% ============================================================================
% IEEE-STYLE JOURNAL ARTICLE
% HTP as Embedding Layer for Micro-LLMs: A Preliminary Study
% ============================================================================
\documentclass[10pt,twocolumn,a4paper]{article}

\usepackage[utf8]{inputenc}
\usepackage[T1]{fontenc}
\usepackage{amsmath,amssymb,amsfonts}
\usepackage{graphicx}
\usepackage{textcomp}
\usepackage{xcolor}
\usepackage{booktabs}
\usepackage{multirow}
\usepackage{hyperref}
\usepackage{cite}
\usepackage{array}
\usepackage{url}
\usepackage{balance}
\usepackage{geometry}
\usepackage{titlesec}
\usepackage{abstract}
\usepackage{enumitem}

\geometry{
    a4paper,
    left=1.7cm, right=1.7cm,
    top=2.0cm, bottom=2.0cm,
    columnsep=0.6cm
}

\hypersetup{
    colorlinks=true,
    linkcolor=blue, citecolor=blue, urlcolor=blue,
    pdftitle={HTP as Embedding Layer for Micro-LLMs: A Preliminary Study},
    pdfauthor={Tcharlies Bachmann Schmitz}
}

\titleformat{\section}{\normalfont\large\bfseries\centering}{\Roman{section}.}{0.5em}{}
\titleformat{\subsection}{\normalfont\normalsize\itshape}{\Alph{subsection}.}{0.5em}{}
\titleformat{\subsubsection}{\normalfont\normalsize\itshape}{\arabic{subsubsection})}{0.5em}{}
\setlist[enumerate]{nosep,leftmargin=*}
\setlist[itemize]{nosep,leftmargin=*}

\newcommand{\htpname}{HTP}
\newcommand{\htpmodel}{HTP-LLM}
\newcommand{\stdmodel}{Std-LLM}

\begin{document}

% ============================================================================
% TITLE
% ============================================================================
\twocolumn[
\begin{@twocolumnfalse}
\begin{center}
\Large\bfseries Harmonic Token Projection as Embedding Layer\\for Micro-Scale Language Models:\\A Preliminary Comparison Against Learned Embeddings
\vspace{0.5cm}

\normalsize
\textit{Tcharlies Bachmann Schmitz}\\
Independent Researcher\\
tcharliesschmitz@gmail.com\\
Brazil\\
\vspace{0.5cm}
\end{center}

\begin{center}
\parbox{0.9\textwidth}{
\textbf{Abstract}---This paper reports a preliminary experiment comparing Harmonic Token Projection (\htpname{}), a deterministic embedding method, against standard learned embeddings in the context of micro-scale language models. Two models sharing the same transformer backbone are trained on a small Portuguese-language corpus (2,862 tokens) and evaluated on 20 manually scored prompts. In this limited setting, the \htpname{}-based model (220,944 parameters) produced correct responses on 9 of 20 prompts, compared to 7 of 20 for the standard model (253,840 parameters). The \htpname{} model also generated fewer incoherent outputs (6 vs.\ 8). We observe that the two models exhibit different failure patterns: the standard model is more prone to syntactic breakdown, while the \htpname{} model tends to drift to related but incorrect topics. These observations, while limited by the small scale of the experiment, suggest that \htpname{} may warrant further investigation as an embedding strategy for parameter-constrained language models. We provide reproduction instructions so that other researchers can verify and challenge these findings.

\vspace{0.3cm}
\textbf{Keywords}---Harmonic Token Projection, Language Models, Parameter-Free Embeddings, Transformer, Reproducibility.
}
\end{center}
\vspace{0.5cm}
\end{@twocolumnfalse}
]

% ============================================================================
% I. INTRODUCTION
% ============================================================================
\section{Introduction}
\label{sec:introduction}

Building a language model with fewer than 250,000 parameters involves difficult trade-offs. One of the most significant is the cost of the embedding layer---the component that maps each word in the vocabulary to a dense vector. With a vocabulary of 736 words and an embedding dimension of 80, this layer requires 58,880 parameters, roughly a quarter of the entire model. For larger vocabularies, the proportion grows quickly: a vocabulary of 4,000 tokens would require 320,000 embedding parameters, exceeding the total budget of many micro-scale architectures.

This observation motivated us to explore whether a parameter-free embedding alternative could be practical at this scale. Harmonic Token Projection (\htpname{}) \cite{schmitz2025htp} is a deterministic method that derives word vectors from the Unicode structure of each word through modular arithmetic and trigonometric functions. It requires no training and no learnable parameters. Originally proposed for structured data applications, it had not previously been tested as an embedding layer for autoregressive language models.

This paper reports a small-scale, controlled experiment comparing \htpname{} against standard learned embeddings. We built two models with the same transformer architecture, trained them on the same data under identical conditions, and compared their outputs on 20 prompts. The results are preliminary and should be interpreted with caution given the small corpus and limited evaluation, but they suggest some patterns that may be worth investigating at larger scales.

% ============================================================================
% II. THE COST OF EMBEDDINGS AT SMALL SCALE
% ============================================================================
\section{The Cost of Embeddings at Small Scale}
\label{sec:problem}

To provide context for the experiment, it is worth examining how the parameter budget is distributed in our two models.

The transformer backbone---two layers of multi-head attention and SwiGLU feed-forward networks---requires approximately 136,000 parameters in both cases. This component is responsible for learning sequential patterns in text, and its capacity is identical across models.

The difference lies in the embedding layer. The standard model uses a learned embedding matrix of $736 \times 80 = 58{,}880$ parameters. Together with the output projection and normalization, the total comes to 253,840 trainable parameters, of which about 53.6\% reside in the transformer.

The \htpname{} model replaces the learned embedding with a deterministic lookup table (zero parameters) and adds a small learned complement of 24 dimensions (17,664 parameters) along with a projection layer (8,320 parameters). The total comes to 220,944 trainable parameters, with 61.6\% in the transformer.

\begin{table}[h]
\centering
\caption{Parameter Distribution}
\label{tab:params}
\begin{tabular}{lrr}
\toprule
\textbf{Component} & \textbf{\htpmodel{}} & \textbf{\stdmodel{}} \\
\midrule
Embedding layer & 17,664 & 58,880 \\
Transformer blocks & 136,160 & 136,160 \\
Output head + norm & 58,960 & 58,960 \\
Input projection & 8,320 & --- \\
\midrule
\textbf{Total} & \textbf{220,944} & \textbf{253,840} \\
\textit{Transformer share} & \textit{61.6\%} & \textit{53.6\%} \\
\bottomrule
\end{tabular}
\end{table}

Whether this 13\% parameter difference and the higher transformer share translate into meaningful quality differences is the question we attempted to explore.

% ============================================================================
% III. HOW HTP WORKS
% ============================================================================
\section{How HTP Works}
\label{sec:htp}

\htpname{} converts any word into a fixed-size vector through three steps.

First, the word is interpreted as a number. Each character has a Unicode codepoint, and these are combined into a single large integer by treating them as digits in base 65,536. Special tokens receive small fixed integers.

Second, this integer is divided by each of the first $K$ prime numbers (where $K = d/2$ for a $d$-dimensional embedding), and only the remainders are kept. Because prime numbers are coprime, the combination of remainders uniquely identifies the original integer with high probability.

Third, each remainder is mapped to a pair of sine-cosine values, producing bounded coordinates between $-1$ and $+1$. Concatenating all pairs yields the final embedding vector.

The result is deterministic (the same word always produces the same vector), and requires no learned parameters. The obvious trade-off is that these embeddings carry no semantic information: words that are similar in meaning are not necessarily close in \htpname{} space.

To partially address this limitation, our \htpname{} model includes a small learned embedding of 24 dimensions, concatenated with the 80-dimensional \htpname{} vector and projected back to 80 dimensions. This small component can learn distributional patterns during training, though its capacity is limited.

% ============================================================================
% IV. EXPERIMENTAL DESIGN
% ============================================================================
\section{Experimental Design}
\label{sec:experiment}

The goal was to compare the two embedding strategies while controlling for as many other variables as possible.

Both models use the same transformer: two decoder blocks with four-head attention, SwiGLU feed-forward networks (hidden dimension 176), Rotary Position Embeddings \cite{su2024roformer}, and RMSNorm \cite{zhang2019root}. Both are trained on the same data: a single Portuguese-language text file about Brazilian geography, containing 2,862 word-level tokens across 736 unique words. The text mixes expository paragraphs with question-answer pairs.

Training is identical: AdamW optimizer, learning rate $3 \times 10^{-4}$, weight decay 0.1, batch size 32, sequence length 128, 15 epochs, gradient clipping at norm 1.0, Xavier uniform initialization, CPU only.

For evaluation, we used 20 prompts in three categories: 10 structured Q\&A prompts matching the training format, 7 free-form sentence beginnings, and 3 direct questions. Generation uses temperature 0.7, top-$k$ sampling with $k=30$, and a maximum of 50 tokens. Each response was manually classified as \textit{Correct}, \textit{Wrong} (coherent but incorrect), or \textit{Incoherent} (broken syntax, loops, or nonsense).

We acknowledge that manual evaluation by the authors introduces bias, and that a single training run per model does not account for stochastic variation. These limitations should be kept in mind when interpreting the results.

% ============================================================================
% V. OBSERVATIONS
% ============================================================================
\section{Observations}
\label{sec:results}

\subsection{Training}

Both models converge to similar final loss values. The standard model shows a slight advantage in early epochs (loss 0.91 at epoch 5 vs.\ 1.13 for \htpname{}), which could reflect the benefit of quickly learning useful word representations. By epoch 10, the difference is negligible, and both models finish near 0.09.

\begin{table}[h]
\centering
\caption{Training Loss at Selected Epochs}
\label{tab:convergence}
\begin{tabular}{ccc}
\toprule
\textbf{Epoch} & \textbf{\htpmodel{}} & \textbf{\stdmodel{}} \\
\midrule
1 & 5.629 & 5.637 \\
5 & 1.128 & 0.908 \\
10 & 0.179 & 0.180 \\
15 & 0.087 & 0.090 \\
\bottomrule
\end{tabular}
\end{table}

It would be premature to draw strong conclusions from the small final difference (0.087 vs.\ 0.090), as both values indicate substantial memorization of the training data.

\subsection{Generation Quality}

On the 20-prompt benchmark, we observed the following distribution:

\begin{table}[h]
\centering
\caption{Benchmark Results (20 Prompts)}
\label{tab:aggregate}
\begin{tabular}{lcc}
\toprule
& \textbf{\htpmodel{}} & \textbf{\stdmodel{}} \\
\midrule
Correct & 9 (45\%) & 7 (35\%) \\
Wrong & 5 (25\%) & 5 (25\%) \\
Incoherent & 6 (30\%) & 8 (40\%) \\
\midrule
Parameters & 220,944 & 253,840 \\
\bottomrule
\end{tabular}
\end{table}

The \htpname{} model produced correct responses on 9 prompts vs.\ 7 for the standard model, and incoherent responses on 6 vs.\ 8. However, we caution that with only 20 prompts and no statistical testing, this difference could partly reflect chance. The sample size is too small for confident claims about one model being ``better'' than the other.

Breaking down by prompt category, we noticed that the \htpname{} model seemed to perform relatively better on free-form completions (4/7 vs.\ 2/7), while the standard model had an advantage on direct questions (1/3 vs.\ 0/3). On structured Q\&A prompts, results were similar (5/10 vs.\ 4/10).

\subsection{Qualitative Differences}

Beyond the aggregate numbers, we observed some patterns in the types of responses each model produces, though we hesitate to generalize from so few examples.

On prompts with near-exact matches in the training data---such as ``Qual \'e o maior bioma do Brasil?''---both models produce essentially identical, correct responses. This is expected: both have memorized the relevant passage.

On prompts requiring the model to compose text without a direct template, differences emerge. When asked to complete ``A regi\~ao nordeste do Brasil,'' the \htpname{} model produced a fluent paragraph about the zona da mata, cana-de-a\c{c}\'ucar, and the agreste. The standard model generated text about the Paran\'a River, which is unrelated to the Northeast. On ``O clima do Brasil,'' the \htpname{} model described regional climate variation; the standard model produced fragmented text.

However, on ``Qual \'e a capital do Brasil'' (a direct question without the Pergunta/Resposta markers), the standard model correctly answered with Bras\'ilia, Kubitschek, and the year 1960, while the \htpname{} model drifted to unrelated content.

\subsection{Failure Patterns}

We noticed a tentative pattern in how the models fail, though we want to be careful not to overstate observations from a small sample.

When the \htpname{} model produces an incorrect response, it often generates coherent text about a related topic---describing the Sert\~ao when asked about the Cerrado, or the Prata basin when asked about the Pantanal. When the standard model fails, it more frequently produces syntactically broken text, including repetitive loops (in one case, repeating ``caatinga?'' five times in a row).

If this pattern were to hold at larger scales, it might suggest that the additional transformer capacity in the \htpname{} model helps maintain syntactic coherence even when factual retrieval fails. But we emphasize that this is a tentative observation from a limited experiment, not a validated finding.

% ============================================================================
% VI. POSSIBLE EXPLANATIONS
% ============================================================================
\section{Possible Explanations}
\label{sec:interpretation}

We offer some speculative explanations for the patterns we observed, while acknowledging that alternative explanations may exist.

The \htpname{} model's relatively stronger performance on free-form completions might be related to the higher proportion of parameters in its transformer backbone (61.6\% vs.\ 53.6\%). If sequential pattern learning---syntax, word order, sentence structure---is primarily the transformer's responsibility, then having more capacity there could help. However, we cannot rule out that other factors, such as the specific initialization or the interaction between \htpname{} representations and the transformer, play a role.

The standard model's advantage on the ``capital'' question might reflect the ability of learned embeddings to develop direct token associations during training. In a learned embedding space, ``capital'' and ``bras\'ilia'' may become proximate during training, creating a shortcut that the \htpname{} model cannot exploit. This is plausible, but a single example is insufficient evidence.

The difference in failure modes---if it is real and not an artifact of the small sample---might reflect the trade-off between embedding capacity and transformer capacity. A model that devotes more parameters to embeddings may have stronger word-level representations but weaker sequential modeling, leading to syntactic breakdown when the model encounters unfamiliar patterns. A model with weaker embeddings but a stronger transformer might maintain syntactic coherence while making factual errors. This is speculative, and we present it as a hypothesis for future testing rather than a conclusion.

% ============================================================================
% VII. SCALING CONSIDERATIONS
% ============================================================================
\section{Scaling Considerations}
\label{sec:scaling}

One aspect of \htpname{} that does not require speculation is its scaling behavior with vocabulary size. The embedding cost of the standard model grows linearly with vocabulary ($V \times d$), while the \htpname{} model's embedding cost grows only through the small learned component ($V \times 24$) and the fixed output head ($V \times 80$). The \htpname{}-specific components (deterministic table and input projection) are independent of vocabulary size.

\begin{table}[h]
\centering
\caption{Projected Parameters by Vocabulary Size}
\label{tab:scaling}
\begin{tabular}{lccc}
\toprule
\textbf{Vocab} & \textbf{\htpmodel{}} & \textbf{\stdmodel{}} & \textbf{Difference} \\
\midrule
736 & 221K & 254K & 13\% \\
2,000 & 251K & 355K & 29\% \\
4,000 & 291K & 515K & 43\% \\
8,000 & 371K & 835K & 56\% \\
\bottomrule
\end{tabular}
\end{table}

At larger vocabularies, the parameter savings become substantial. Whether these savings translate into quality improvements---or whether learned embeddings become increasingly important as vocabulary grows---is an empirical question that this paper does not answer.

% ============================================================================
% VIII. LIMITATIONS
% ============================================================================
\section{Limitations}
\label{sec:limitations}

This experiment has significant limitations that constrain the strength of any conclusions.

The training corpus contains only 2,862 tokens. This is orders of magnitude smaller than typical language model training sets. At this scale, both models are likely memorizing rather than learning to generalize. Results could be very different with larger training data.

The evaluation is based on 20 prompts manually scored by the authors. This introduces subjectivity, and the sample size is too small for statistical significance testing. A larger evaluation with multiple annotators and automated metrics would be necessary to make confident claims.

All data is in Portuguese and concerns a single domain (geography). We do not know how \htpname{} would perform in other languages, particularly those with different scripts, or across multiple domains.

Each model was trained once. Without multiple runs with different random seeds, we cannot separate the effect of the embedding strategy from the effect of initialization luck.

\htpname{} embeddings carry no semantic information. They encode the character sequence of a word, not its meaning. In tasks or at scales where semantic similarity in the embedding space is important, this limitation could be significant.

Finally, our comparison is limited to a single model size and configuration. The relative merits of \htpname{} and learned embeddings might change with different model dimensions, depths, or training regimes.

% ============================================================================
% IX. REPRODUCING THIS WORK
% ============================================================================
\section{Reproducing This Work}
\label{sec:reproducibility}

Given the preliminary nature of this work, independent reproduction is especially important. We have aimed to make this as straightforward as possible.

All source code, training data, benchmark prompts, and this paper are available in a public repository \cite{htpbenchmark2025}. The experiment requires Python 3.10 or later and PyTorch. No GPU is needed. A single script (\texttt{run\_benchmark.py}) trains both models under identical conditions, generates responses for all 20 prompts, and prints a side-by-side comparison. The full benchmark runs in approximately 30 minutes on a modern CPU.

Training uses AdamW with learning rate $3 \times 10^{-4}$, weight decay 0.1, batch size 32, sequence length 128, 15 epochs, and gradient clipping at norm 1.0. The training data is a single text file about Brazilian geography. The 20-prompt benchmark uses temperature 0.7, top-$k$ sampling with $k = 30$, and 50 maximum tokens.

Due to stochastic training, exact outputs will vary between runs. We encourage researchers not only to reproduce the experiment but to improve upon it: using larger corpora, automated evaluation metrics, multiple seeds, and cross-domain testing. If \htpname{}'s advantages disappear under more rigorous conditions, that is a valuable finding in itself.

% ============================================================================
% X. CONCLUSION
% ============================================================================
\section{Conclusion}
\label{sec:conclusion}

This paper reports a small, controlled experiment comparing two embedding strategies for micro-scale language models: Harmonic Token Projection, a deterministic method requiring no learnable parameters, and the standard learned embedding used in virtually all neural language models.

In our limited setting, the \htpname{}-based model produced correct responses on 9 of 20 prompts compared to 7 for the standard model, and fewer incoherent outputs (6 vs.\ 8), while using 13\% fewer parameters. We also observed tentative differences in failure patterns, with the \htpname{} model maintaining greater syntactic coherence when producing incorrect answers.

We do not claim that these results demonstrate \htpname{}'s superiority. The experiment is too small, the evaluation too limited, and the possible confounding factors too numerous for such a claim. What we do suggest is that \htpname{} appears to be a viable embedding strategy at this scale, and that the patterns we observed---particularly the relationship between parameter allocation and output coherence---may be worth investigating further.

The embedding layer is so deeply assumed in language model design that alternatives are rarely considered. This paper is an invitation to consider one.

% ============================================================================
% REFERENCES
% ============================================================================
\bibliographystyle{ieeetr}
\bibliography{references}

\balance

\end{document}
